\section{Algebraic Structures}

An algebraic structure is nothing but a set (or maybe a collection of sets) bundled with one or more operations. The study 
of algebraic structures is itself the study of the properties of the operations defined over the element of the set. Please keep in 
mind that the same set can be used to define multiple algebraic structures. For example, the set of integers can be used to define
a group, a ring, a field, etc. The operations defined over the set are the key to defining the algebraic structure.

\subsection{Operations}
"Operation" is just a fancy term to refer functions in the context of algebraic structures. An operation can be either 
internal or external. An internal operation is a function that takes some elements from the set and returns an element from the set. Internal 
operations are the most common type of operations in algebraic structures.

\subsubsection{Basics of operations and special elements}
In order to properly study operations in algebraic structures, we need to give a formal definition of an operation and 
its properties. We will also introduce the concept of a special element for a given operation, which is an element of the algebraic 
structure that has some special properties with respect to the operation at hand.

\begin{definition}
    An internal binary operation $\perp$ over a set $S$ is a function $\perp: S \times S \rightarrow S$
\end{definition}

\begin{definition}
    An external binary operation $\cdot$ over the pair of sets $K$, $V$ is a function defined as $\cdot: K \times V \rightarrow V$, note 
    that this definition encompasses the case where $K = V$ hence its a generalization of the internal binary operation.
\end{definition}

\begin{definition}
    An unary operation $!$ over a set $S$ is a function $!: S \rightarrow S$, note that 
    an unary operation can never be external so specifying that it is internal is redundant.
\end{definition}

\begin{definition}
    With respect to a binary internal operation $\perp$ over a set $S$, an element $\epsilon \in S$ is said to be the left-identity-element of 
    the operation $\perp$ if $\forall x \in S: \epsilon \perp x = x$.
\end{definition}

\begin{definition}
    With respect to a binary internal operation $\perp$ over a set $S$, an element $\delta \in S$ is said to be the right-identity-element of 
    the operation $\perp$ if $\forall x \in S: x \perp \delta = x$.
\end{definition}

Given these very few definitions we can already prove some interesting statements about the properties of the identity elements of an operation. In 
particular, we can prove that if a set has both a left-identity-element and a right-identity-element, then they are the are both simultaneously the
left-identity-element and the right-identity-element of the operation. We can also prove that under these circumstances the two elements are not distinct, meaning 
that indeed there can only be one identity element for a given operation both on the left and on the right. This property is very important in the study of algebraic
structures, as it allows us to define the concept of identity element in a more general way.

\newpage

\begin{lemma}
    If two elements $\epsilon, \delta \in S$ exist such that they are left-identity-element and right-identity-element of the operation 
    respectively, then $\epsilon = \delta$.
\end{lemma}
\begin{proof}
    Let $\epsilon \in S$ be the left-identity-element $\perp$, then let $\delta \in S$ be the right-identity-element of $\perp$. The following statements 
    are then true by the very definition of left and right identity elements:

    \begin{align*}
        \forall x \in S: \epsilon \perp x = x \\
        \forall x \in S: x \perp \delta = x
    \end{align*}

    By the definition of the right-identity-element, we can substitute $x = \epsilon$. By the definition of the 
    left-identity-element, we can also substitute $x = \delta$. Ending up with:

    \begin{align*}
        \epsilon \perp \delta = \epsilon \\
        \epsilon \perp \delta = \delta
    \end{align*}

    By the transitive property of equality, we can then conclude that $\epsilon = \delta$. Note that this means that 
    the left-identity-element and the right-identity-element are the same element, and that the element itslef is therefore 
    the identity element of the operation $\perp$ on both sides (this holds true regardless of wether $\perp$ has any other property or not).

    \begin{align*}
        \epsilon = \epsilon \perp \delta = \delta
    \end{align*}

\end{proof}


\begin{definition}
    With respect to a binary internal operation $\perp$ over a set $S$, an element $\hat{a}$ is said to be the left-inverse-element of
    an element $a \in S$ if $\hat{a} \perp a = \epsilon$ where $\epsilon$ is the left-identity-element of the operation $\perp$.
\end{definition}

\begin{definition}
    With respect to a binary internal operation $\perp$ over a set $S$, an element $\bar{a}$ is said to be the right-inverse-element of
    an element $a \in S$ if $a \perp \bar{a} = \delta$ where $\delta$ is the right-identity-element of the operation $\perp$.
\end{definition}

With these definitions we can now go ahead and prove an interesting property of the inverse elements of an operation. In particular, we can prove that
if an element has both a left-inverse-element and a right-inverse-element, then the two elements are the same Under the assumption that the operation is 
associative, we can also prove that the inverse element is unique. \\

Please note how this also automatically implies that the identity element is unique since 
the very existence of both a left and a right inverse element implies that there exists both a left and a right identity element, which itself implies that 
the two elements are the same and therefore unique. 